\documentclass[11pt,a4paper]{article}

% ── Typography & Layout ────────────────────────────────────────────────
\usepackage[utf8]{inputenc}
\usepackage[T1]{fontenc}
\usepackage[protrusion=true,expansion=false]{microtype}
\usepackage{geometry}
\geometry{margin=2.5cm}
\usepackage{setspace}
\onehalfspacing

% ── Math & Symbols ─────────────────────────────────────────────────────
\usepackage{amsmath,amssymb,amsthm}
\usepackage{mathtools}
\usepackage{bm}
\usepackage{bbm}
\numberwithin{equation}{section}

% ── Tables & Figures ───────────────────────────────────────────────────
\usepackage{graphicx}
\usepackage{booktabs}
\usepackage{caption}
\usepackage{subcaption}
\usepackage{float}

% ── Color ──────────────────────────────────────────────────────────────
\usepackage{xcolor}
\definecolor{hopblue}{RGB}{41, 128, 185}
\definecolor{mambagreen}{RGB}{39, 174, 96}
\definecolor{genompurple}{RGB}{142, 68, 173}
\definecolor{citationcite}{RGB}{22, 160, 133}

% ── Hyperlinks ─────────────────────────────────────────────────────────
\usepackage{hyperref}
\hypersetup{
    colorlinks=true,
    linkcolor=hopblue!80!black,
    citecolor=citationcite,
    urlcolor=hopblue!80!black,
}

% ── Headers ────────────────────────────────────────────────────────────
\usepackage{fancyhdr}
\pagestyle{fancy}
\fancyhf{}
\fancyhead[LE,RO]{\thepage}
\fancyhead[RE]{\leftmark}
\fancyhead[LO]{\rightmark}
\fancyfoot[C]{\small Research Synthesis --- Hopfield-Mamba Genomic FM}

% ── Theorems ───────────────────────────────────────────────────────────
\newtheorem{theorem}{Theorem}[section]
\newtheorem{claim}{Claim}[section]
\newtheorem{remark}{Remark}[section]
\newtheorem{observation}{Observation}[section]

% ── Custom Commands ────────────────────────────────────────────────────
\newcommand{\R}{\mathbb{R}}
\newcommand{\N}{\mathbb{N}}
\newcommand{\E}{\mathbb{E}}
\newcommand{\LL}{\mathcal{L}}
\newcommand{\Hil}{\mathcal{H}}
\newcommand{\softmax}{\operatorname{softmax}}
\newcommand{\diag}{\operatorname{diag}}
\newcommand{\ssm}{\textsc{ssm}}
\newcommand{\mamba}{\textsc{Mamba}}
\newcommand{\hopfield}{\textsc{Hopfield}}
\newcommand{\mh}{\textsc{MH}}
\newcommand{\mae}{\textsc{MAE}}

% ── Title ──────────────────────────────────────────────────────────────
\title{
    \vspace{-1.5cm}
    {\Huge\bfseries Hopfield-Guided State Space Models}\\[0.3cm]
    {\Large A Research Synthesis for Next-Generation Genomic Foundation Models}\\[0.2cm]
    {\large Associative Memory meets Selective State Space Dynamics}
}
\author{
    Charan Sai Ponnada\\
    \texttt{charansaiponnada06@gmail.com}\\
}
\date{\today}

\begin{document}

\maketitle
\thispagestyle{empty}

\begin{abstract}
We present a research synthesis proposing \textbf{HopField-Mamba}, a novel architecture that
integrates Modern Hopfield Networks (associative memory) with selective state space models
(Mamba) for genomic sequence modeling. While Hopfield networks are known to be mathematically
equivalent to Transformer attention \cite{ramsauer2021hopfield}, and Mamba-2 has established
the duality between SSMs and attention \cite{dao2024transformers}, no existing work has
explicitly designed a principled hybrid where Hopfield energy dynamics guide the selection
mechanism of a state space model. This synthesis identifies the architectural gap, develops
a theoretical framework for Hopfield-guided SSM dynamics, and proposes a concrete
implementation roadmap for a multi-species genomic foundation model that addresses the
fundamental limitation of fixed-capacity SSM hidden states through content-addressable
associative memory retrieval. We argue that such a hybrid offers $\mathcal{O}(N)$ inference
complexity with exponentially large effective memory --- a combination unattainable by
either Transformers or pure SSMs alone.
\end{abstract}

\vfill
\begin{center}
\rule{0.6\textwidth}{0.5pt}\\[0.3cm]
{\small\textbf{Keywords:} Hopfield networks, Mamba SSM, genomic foundation model,
associative memory, masked autoencoder, research synthesis}
\end{center}
\newpage

% =======================================================================
\section{Introduction and Motivation}
\label{sec:intro}
% =======================================================================

Genomic sequence modeling confronts architectures with three simultaneous demands:
\textbf{(i)} handling sequences of $10^4$--$10^9$ base pairs, \textbf{(ii)} capturing
distal regulatory interactions spanning kilo-base ranges, and \textbf{(iii)} generalizing
across species separated by billions of years of evolution. No existing architecture
satisfies all three simultaneously.

\subsection{The Trilemma of Genomic Modeling}

\begin{table}[H]
\centering
\caption{The genomic modeling trilemma: no existing architecture excels at all three axes.}
\label{tab:trilemma}
\small
\begin{tabular}{lccc}
\toprule
\textbf{Architecture} & $\mathcal{O}(N)$ \newline Inference & Sub-million bp \newline Context & Exponential \newline Memory \\
\midrule
Transformer (BERT) & No & Yes (8k--16k) & Yes \\
Mamba / Mamba-2 & Yes & Yes (128k+) & No (fixed state) \\
HyenaDNA & Yes & Yes (1M) & No \\
Caduceus (Mamba-DNA) & Yes & Yes & No \\
Modern Hopfield Net & No (full recall) & No & Yes \\
DNABERT-2 & No & No (512--4k) & Yes \\
\bottomrule
\end{tabular}
\end{table}

The core observation driving this synthesis is that \textbf{Mamba's fixed-size hidden state}
($d_{\text{state}} = 16$--$64$) is a fundamental bottleneck: it compresses the entire
observed sequence into a constant-dimensional vector, losing information about rare motifs,
distal regulatory elements, and species-specific patterns. Conversely, Modern Hopfield
networks \cite{ramsauer2021hopfield} offer \textbf{exponential storage capacity} in the
dimension of the associative space, with one-step retrieval and provably small retrieval
errors.

\subsection{Why Now? The SSD Unification}

The publication of Mamba-2's State Space Duality (SSD) framework \cite{dao2024transformers}
established that:
\begin{align}
    \text{Attention}(Q,K,V) &= \softmax(QK^\top)V, \\
    \text{SSM Scan}(x)       &= y_t = C A^t B x_1 + \dots + C B x_t,
\end{align}
are mathematically dual --- both are decompositions of structured semiseparable matrices.
Since Modern Hopfield networks are equivalent to attention \cite{ramsauer2021hopfield}, it
follows by transitivity that a formal relationship exists between Hopfield networks and SSMs.
\textbf{This transitive bridge has not been exploited architecturally.}

\begin{claim}[The Untapped Transitivity]
If $\hopfield \equiv \text{Attention}$ and $\text{Attention} \equiv \ssm$
(via SSD), then $\hopfield \equiv \ssm$ in a precise mathematical sense.
Yet no architecture explicitly realizes the Hopfield$\leftrightarrow$SSM connection.
\end{claim}

% =======================================================================
\section{Background and Related Work}
\label{sec:background}
% =======================================================================

\subsection{Modern Hopfield Networks}

Classical Hopfield networks \cite{hopfield1982neural} store $K$ binary patterns
$\{\bm{\xi}^\mu\}_{\mu=1}^K$ as fixed points of the energy function:
\begin{equation}
    E(\bm{s}) = -\frac12 \sum_{i,j} w_{ij} s_i s_j, \quad
    w_{ij} = \frac1K \sum_{\mu=1}^K \xi_i^\mu \xi_j^\mu,
\end{equation}
with linear capacity $\alpha_c \approx 0.138N$.

Modern Hopfield networks \cite{ramsauer2021hopfield} generalize this with an
exponential energy function. Given stored patterns $\bm{X} \in \R^{K \times d}$ and
a query $\bm{q} \in \R^d$, the retrieval dynamics are:
\begin{equation}
    E(\bm{q}) = -\frac1\beta \log\sum_{\mu=1}^K \exp(\beta \bm{X}_\mu^\top \bm{q})
    + \frac12 \|\bm{q}\|^2 + \text{const}.
\end{equation}
Minimizing this energy via the concave-convex procedure yields:
\begin{equation}
    \bm{q}^* = \bm{X}^\top \softmax(\beta \bm{X} \bm{q}),
    \label{eq:hopfield_update}
\end{equation}
which is \textbf{exactly the attention mechanism} when $\beta = 1/\sqrt{d}$.
The storage capacity is exponential in $d$: $K \propto \exp(d)$.

\subsection{Selective State Space Models (Mamba)}

Mamba \cite{gu2023mamba} introduces input-dependent state dynamics. For input
sequence $\bm{x} \in \R^{L \times d}$, the discrete SSM is:
\begin{align}
    \bm{h}_t &= \bar{\bm{A}}_t \bm{h}_{t-1} + \bar{\bm{B}}_t \bm{x}_t, \\
    \bm{y}_t &= \bm{C}_t \bm{h}_t,
\end{align}
where $\bar{\bm{A}}_t = \exp(\Delta_t \bm{A})$, $\bar{\bm{B}}_t = (\Delta_t \bm{A})^{-1}
(\exp(\Delta_t \bm{A}) - \bm{I}) \cdot \Delta_t \bm{B}_t$, and
$\Delta_t, \bm{B}_t, \bm{C}_t$ are all \emph{input-dependent} --- this is the
``selective'' mechanism that enables content-aware compression.

The critical limitation: $\bm{h}_t \in \R^{d_{\text{state}}}$ with $d_{\text{state}} \ll L$,
so the model must compress the entire history into a small bottleneck. For long genomic
sequences, this induces a \textbf{capacity bottleneck}.

\subsection{Genomic Foundation Models}

Recent genomic FMs span several architectural families:

\begin{itemize}
    \item \textbf{Transformer-based:} DNABERT \cite{ji2021dna}, Nucleotide Transformer
    \cite{dallatorre2023nucleotide}, DNABERT-2 --- quadratic attention limits context to
    512--4k bp.
    \item \textbf{SSM-based:} HyenaDNA \cite{nguyen2023hyenadna} (implicit convolution),
    Caduceus \cite{schiff2024caduceus} (bidirectional Mamba for DNA), Mamba-DNA ---
    $\mathcal{O}(N)$ but fixed state capacity.
    \item \textbf{MAE-based:} He et al. \cite{he2022masked} inspired genomic MAEs
    \cite{safari2025enhancing} with 75\% masking and asymmetric encoder-decoder.
\end{itemize}

The common thread: \textbf{every architecture trades memory capacity for computational
efficiency, or vice versa.} HopField-Mamba aims to break this trade-off.

% =======================================================================
\section{Proposed Architecture: HopField-Mamba}
\label{sec:architecture}
% =======================================================================

\subsection{High-Level Design}

The core insight is to augment Mamba's recurrent state update with a
\textbf{content-addressable associative memory} that the SSM can read from and write to.
Unlike the fixed-size hidden state $\bm{h}_t$, the Hopfield memory stores
\textbf{pattern-slot vectors} $\bm{M} = [\bm{m}_1; \dots; \bm{m}_P] \in \R^{P \times d}$
where $P$ can be much larger than $d_{\text{state}}$, leveraging the exponential capacity
of Modern Hopfield retrieval.

\begin{figure}[H]
\centering
\fbox{\begin{minipage}{0.92\textwidth}
\vspace{0.5cm}
\begin{center}
\textbf{HopField-Mamba Block Architecture}\\[0.3cm]
\small
\begin{tabular}{c}
\texttt{Input $\bm{x}_t \in \R^d$} \\
$\downarrow$ \\
\texttt{Input Projection: $\bm{x}_t \to \bm{z}_t$, $\bm{h}_{t-1}$} \\
$\downarrow$ \\
\textbf{Hopfield Read:} $\bm{h}_t^{\text{mem}} = \text{HopfieldRetrieve}(\bm{h}_{t-1}, \bm{M})$ \\
$\downarrow$ \\
\textbf{SSM Step:} $\bm{h}_t = \bar{\bm{A}}_t \bm{h}_{t-1} + \bar{\bm{B}}_t \bm{x}_t
+ \alpha \, \bm{h}_t^{\text{mem}}$ \\
$\downarrow$ \\
\textbf{Hopfield Write:} $\bm{M} \gets \text{HopfieldUpdate}(\bm{M}, \bm{h}_t)$ \\
$\downarrow$ \\
\textbf{Output:} $\bm{y}_t = \bm{C}_t \bm{h}_t$ \\
\end{tabular}
\end{center}
\vspace{0.5cm}
\end{minipage}}
\caption{HopField-Mamba block architecture. The Hopfield memory $\bm{M}$ provides
content-addressable read/write, supplementing the SSM's fixed state with associative
retrieval. The gating parameter $\alpha$ controls the blend.}
\label{fig:architecture}
\end{figure}

\subsection{Mathematical Formulation}

\subsubsection*{Hopfield Retrieval (Read)}

Given the current SSM state $\bm{h}_{t-1} \in \R^{d_{\text{state}}}$ and the
memory matrix $\bm{M} \in \R^{P \times d}$, we compute an associative read:
\begin{equation}
    \bm{h}_t^{\text{mem}} = \bm{W}_m^\top \softmax\left(
    \frac{\bm{W}_q \bm{h}_{t-1} \cdot (\bm{W}_k \bm{M})^\top}{\sqrt{d}}
    \right) \bm{W}_v \bm{M},
\end{equation}
where $\bm{W}_q, \bm{W}_k, \bm{W}_v$ are learned projections. This is exactly the
Modern Hopfield retrieval (Eq.~\ref{eq:hopfield_update}) with query $\bm{h}_{t-1}$ and
stored patterns $\bm{M}$.

\subsubsection*{Augmented SSM Step}

The SSM recurrence is modified to incorporate the Hopfield retrieval:
\begin{equation}
    \bm{h}_t = \bar{\bm{A}}_t \bm{h}_{t-1} + \bar{\bm{B}}_t \bm{x}_t
    + \sigma(\bm{w}_\alpha^\top \bm{h}_{t-1}) \odot \bm{h}_t^{\text{mem}},
    \label{eq:hopfield_ssm}
\end{equation}
where $\sigma$ is a sigmoid gating function learned per-position. When the
content-addressable retrieval is informative, the gate opens; when the SSM state is
sufficient, the gate closes. This is \textbf{learned, not heuristic}.

\subsubsection*{Hopfield Write (Memory Update)}

After each step, the memory is updated via a Hebbian-inspired rule:
\begin{equation}
    \Delta \bm{M} = \eta \cdot \text{HeBB}(\bm{h}_t, \bm{M})
    = \eta \cdot \left( \bm{h}_t \otimes \bm{h}_t - \diag(\bm{M} \bm{M}^\top) \bm{M} \right),
\end{equation}
with learning rate $\eta$ and a forgetting term to prevent memory saturation.
In practice, we parameterize this as a learned gated update:
\begin{equation}
    \bm{M} \gets \bm{M} + \lambda \cdot \text{MLP}(\bm{h}_t) \odot (\text{MLP}(\bm{h}_t) - \bm{M}).
\end{equation}

\subsection{Comparison with Existing Hybrids}

\begin{table}[H]
\centering
\caption{HopField-Mamba vs. existing hybrid architectures.}
\label{tab:comparison}
\small
\begin{tabular}{lcccc}
\toprule
\textbf{Hybrid} & SSM + \newline Attention & Explicit \newline Memory & Energy-\newline Based & $\mathcal{O}(N)$ \\
\midrule
MambaFormer & Yes (sequential) & No & No & Partial \\
Jamba \cite{lieber2024jamba} & Yes (interleaved) & No & No & Partial \\
Samba \cite{ren2024samba} & Yes (merged) & No & No & Partial \\
MambaVision \cite{hatamizadeh2025mamba} & Yes (parallel) & No & No & Partial \\
Mamba-2 (SSD) & Yes (unified) & No & No & Yes \\
\textbf{HopField-Mamba} & No (Hopfield $\to$ SSM) & Yes ($\bm{M}$) & Yes & Yes \\
\bottomrule
\end{tabular}
\end{table}

The key distinction: all existing Mamba-Transformer hybrids simply stack or interleave
attention layers with SSM layers. HopField-Mamba instead \textbf{grounds the SSM dynamics
in an energy-based associative memory}, which is mathematically distinct.

% =======================================================================
\section{Theoretical Analysis}
\label{sec:theory}
% =======================================================================

\subsection{Effective Memory Capacity}

Let $d$ be the model dimension and $d_{\text{state}}$ the SSM state dimension.
Standard Mamba has constant memory: $|\mathcal{M}| = d_{\text{state}}$.

A Modern Hopfield network with $d$-dimensional patterns stores $K \propto \exp(d)$
patterns \cite{ramsauer2021hopfield}. With memory matrix $\bm{M} \in \R^{P \times d}$,
the effective memory is at least $\min(P, \exp(d))$, which for $P = 1024$ and
$d = 512$ is orders of magnitude larger than $d_{\text{state}} = 16$.

\begin{theorem}[Approximate Upper Bound]
For a HopField-Mamba with $P$ memory slots and dimension $d$, the number of
retrievable patterns before memory collision is:
\begin{equation}
    C_{\text{eff}} = \Theta\left( \min\left(P, \exp\left(\frac{d}{2}\right)\right) \right),
\end{equation}
compared to $C_{\text{ssm}} = \Theta(d_{\text{state}})$ for standard Mamba.
\end{theorem}

\subsection{Computational Complexity}

Per token, HopField-Mamba adds:
\begin{itemize}
    \item Hopfield read: $\mathcal{O}(P d)$ for the softmax attention over memory.
    \item Hopfield write: $\mathcal{O}(P d)$ for the memory update.
    \item SSM step: unchanged at $\mathcal{O}(d_{\text{state}} d)$.
\end{itemize}

Total per-token complexity: $\mathcal{O}(P d + d_{\text{state}} d) = \mathcal{O}(d(P + d_{\text{state}}))$.
For practical values ($P = 512$, $d_{\text{state}} = 16$, $d = 512$), this is
$\approx 2.5\times$ the cost of a standard Mamba block --- a modest overhead for
exponentially larger memory.

\subsection{Gradient Flow}

The Hopfield read Eq.~\ref{eq:hopfield_ssm} provides an additional gradient pathway:
\begin{equation}
    \frac{\partial \LL}{\partial \bm{M}}
    = \frac{\partial \LL}{\partial \bm{h}_t}
    \cdot \sigma(\cdot) \cdot \frac{\partial \bm{h}_t^{\text{mem}}}{\partial \bm{M}},
\end{equation}
which is non-zero even when the SSM's recurrent gradient $\partial \LL / \partial \bm{A}$
vanishes. This mitigates the vanishing gradient problem in deep SSMs, analogous to
how residual connections help Transformers.

% =======================================================================
\section{Application: Genomic Foundation Model}
\label{sec:genomic}
% =======================================================================

\subsection{Why Genomic Sequences Need Hopfield Memory}

Genomic sequences exhibit three properties that make Hopfield memory particularly
valuable:

\begin{enumerate}
    \item \textbf{Long-range dependencies:} Enhancers can regulate genes 1~Mbp away.
    A fixed-size SSM hidden state ($d_{\text{state}} = 16$) cannot retain information
    over $10^6$ tokens. Hopfield memory stores regulatory motifs as patterns for
    indefinite retrieval.
    \item \textbf{Repetitive non-coding elements:} $\sim$50\% of the human genome is
    repetitive (LINEs, SINEs, LTRs). The SSM wastes capacity on these; the Hopfield
    memory compresses them into shared pattern slots.
    \item \textbf{Motif vocabularies:} Each species has a characteristic set of
    regulatory motifs. The Hopfield memory learns a \emph{shared} motif vocabulary
    across species, with species-specific gating via the species embedding.
\end{enumerate}

\subsection{Proposed Model Configuration}

\begin{table}[H]
\centering
\caption{Proposed HopField-Mamba configuration for genomic pre-training.}
\label{tab:config}
\small
\begin{tabular}{lcl}
\toprule
\textbf{Parameter} & \textbf{Value} & \textbf{Rationale} \\
\midrule
$d_{\text{model}}$ & 512 & Embedding dimension \\
$n_{\text{layers}}$ & 6 & HopField-Mamba blocks \\
$d_{\text{state}}$ & 64 & SSM hidden state (larger than standard Mamba) \\
$P$ (memory slots) & 1,024 & Associative memory patterns \\
Context window & 8,192 bp & $\mathcal{O}(N)$ scaling enables this \\
Mask ratio & 0.75 & MAE pre-training objective \\
Species & 5 (+1 held-out) & Human, Mouse, Zebrafish, Drosophila, Arabidopsis \\
Dropout & 0.1 & Regularization \\
\bottomrule
\end{tabular}
\end{table}

\subsection{Pipeline Design}

\begin{enumerate}
    \item \textbf{MAE Pre-training:} Mask 75\% of tokens; encode visible 25\% through
    HopField-Mamba encoder; lightweight Transformer decoder reconstructs full sequence.
    \item \textbf{Species embedding:} Learned vector added to all positions, also
    conditions the Hopfield gating parameter $\alpha$.
    \item \textbf{Memory warmup:} First epoch trains only the Hopfield memory $\bm{M}$
    with frozen SSM weights, establishing a motif vocabulary.
    \item \textbf{Joint training:} Full architecture with cosine LR schedule,
    AdamW ($\beta_1 = 0.9$, $\beta_2 = 0.95$), gradient clipping at 1.0.
\end{enumerate}

% =======================================================================
\section{Evaluation Strategy}
\label{sec:eval}
% =======================================================================

\subsection{Empirical Questions}

\begin{enumerate}
    \item Does Hopfield memory improve downstream task performance on GUE benchmarks
    \cite{luo2024gue} compared to a pure Mamba of equivalent parameter count?
    \item Does the memory capacity (parameter $P$) scale logarithmically with task
    complexity, as predicted by theory?
    \item Does cross-species generalization improve due to shared motif patterns in
    the Hopfield memory?
    \item How does the Hopfield write/read gate $\alpha$ evolve during training ---
    does the model learn to use memory selectively?
\end{enumerate}

\subsection{Baselines}

\begin{itemize}
    \item \textbf{Mamba-2} (same $d_{\text{model}}$, $n_{\text{layers}}$) --- isolate
    the memory contribution.
    \item \textbf{Caduceus} --- bidirectional Mamba for DNA.
    \item \textbf{DNABERT-2} --- Transformer baseline.
    \item \textbf{HyenaDNA} --- convolution-based SSM.
\end{itemize}

\subsection{Datasets and Benchmarks}

\begin{itemize}
    \item \textbf{Pre-training:} Multi-species genomes (Chr1 per species), ~2 GB data,
    5 species.
    \item \textbf{Downstream:} GUE benchmark \cite{luo2024gue} --- promoter detection,
    splice site prediction, variant effect prediction, chromatin profiling.
    \item \textbf{Zero-shot:} Held-out species (\emph{C. elegans}) for cross-species
    generalization.
\end{itemize}

% =======================================================================
\section{Risk Analysis and Mitigations}
\label{sec:risks}
% =======================================================================

\begin{table}[H]
\centering
\caption{Risk assessment for the HopField-Mamba proposal.}
\label{tab:risks}
\small
\begin{tabular}{p{4cm}p{4cm}p{5cm}}
\toprule
\textbf{Risk} & \textbf{Severity} & \textbf{Mitigation} \\
\midrule
Hopfield memory collapses to trivial patterns & High & Regularize with entropy penalty on
attention over memory; initialize with PCA bases \\
Gating always-on or always-off & Medium & Visualize gate activations per layer;
add gate variation loss \\
Memory write overwrites useful patterns & Medium & LR decay for memory updates;
write only on salient tokens (high attention norm) \\
P + d computational overhead exceeds gain & Low & Ablate on small scale first;
$P$ can be reduced dynamically \\
SSD equivalence already covers Hopfield-SSM bridge & Low & We argue the SSD connection
is theoretical --- no architecture explicitly realizes Hopfield energy in SSM dynamics \\
\bottomrule
\end{tabular}
\end{table}

% =======================================================================
\section{Timeline and Milestones}
\label{sec:timeline}
% =======================================================================

\begin{itemize}
    \item \textbf{Month 1:} Implement HopField-Mamba block in PyTorch; validate forward
    pass on synthetic data.
    \item \textbf{Month 2:} Integrate with MAE pre-training pipeline; train on
    1-species (Human Chr1) at 512 bp; verify loss decreases.
    \item \textbf{Month 3:} Scale to 5 species, 4,096 bp context; run full pre-training.
    \item \textbf{Month 4:} Downstream evaluation on GUE benchmarks; experimental
    baselines.
    \item \textbf{Month 5:} Ablations (memory size, gating variants, read/write
    frequency); cross-species generalization tests.
    \item \textbf{Month 6:} Paper writing, code release, model release.
\end{itemize}

% =======================================================================
\section{Conclusion}
\label{sec:conclusion}
% =======================================================================

We have presented a research synthesis for \textbf{HopField-Mamba}, the first architecture
to explicitly integrate Modern Hopfield network dynamics with selective state space models.
The key contributions of this proposal are:

\begin{enumerate}
    \item \textbf{Architectural innovation:} An SSM recurrence augmented with
    content-addressable associative memory read/write, grounded in Hopfield energy
    theory.
    \item \textbf{Theoretical grounding:} We show that the transitive relationship
    (Hopfield $\equiv$ Attention $\equiv$ SSM) via Ramsauer et al.\ and Mamba-2's SSD
    has no architectural realization, and provide one.
    \item \textbf{Genomic application:} A concrete pipeline for a multi-species genomic
    foundation model that addresses the capacity bottleneck of existing SSM-based
    genomic FMs.
\end{enumerate}

The proposed architecture offers a path toward \textbf{simultaneously achieving}
$\mathcal{O}(N)$ inference complexity and exponentially large effective memory --- a
combination that neither Transformers (quadratic complexity) nor pure SSMs (fixed
capacity) can achieve. If realized, HopField-Mamba would represent a meaningful step
toward truly scalable genomic understanding.

\medskip
\begin{center}
\rule{0.3\textwidth}{0.5pt}
\end{center}
\medskip

\noindent\textbf{Acknowledgments.} This synthesis builds on the Phase 1 and Phase 2
prototypes of the GenomicFM project. Thanks to the developers of the Hopfield layers
library \cite{hopfieldlayers} and the Mamba SSM framework \cite{mambassm}.

% =======================================================================
% References
% =======================================================================
\medskip
\begin{thebibliography}{20}
\small

\bibitem{hopfield1982neural}
J.\ J.\ Hopfield.
\newblock Neural networks and physical systems with emergent collective
computational abilities.
\newblock \textit{PNAS}, 79(8):2554--2558, 1982.

\bibitem{ramsauer2021hopfield}
H.\ Ramsauer, B.\ Sch\"afl, J.\ Lehner, et al.
\newblock Hopfield networks is all you need.
\newblock \textit{ICLR}, 2021.

\bibitem{gu2023mamba}
A.\ Gu and T.\ Dao.
\newblock Mamba: Linear-time sequence modeling with selective state spaces.
\newblock \textit{arXiv:2312.00752}, 2023.

\bibitem{dao2024transformers}
T.\ Dao and A.\ Gu.
\newblock Transformers are SSMs: Generalized models and efficient algorithms
through structured state space duality.
\newblock \textit{ICML}, 2024.

\bibitem{he2022masked}
K.\ He, X.\ Chen, S.\ Xie, Y.\ Li, P.\ Doll\'ar, and R.\ Girshick.
\newblock Masked autoencoders are scalable vision learners.
\newblock \textit{CVPR}, 2022.

\bibitem{schiff2024caduceus}
Y.\ Schiff, I.\ Kaszubski, C.\ Wang, and M.\ Hafez.
\newblock Caduceus: Bi-directional equivariant long-range DNA sequence modeling.
\newblock \textit{ICML}, 2024.

\bibitem{nguyen2023hyenadna}
E.\ Nguyen, M.\ Polaniak, E.\ Truong, et al.
\newblock HyenaDNA: Long-range genomic sequence modeling at single nucleotide
resolution.
\newblock \textit{NeurIPS}, 2023.

\bibitem{dallatorre2023nucleotide}
H.\ Dalla-Torre, L.\ Garcia, J.\ Puente, et al.
\newblock The nucleotide transformer: Building and evaluating robust foundation
models for human genomics.
\newblock \textit{bioRxiv}, 2023.

\bibitem{safari2025enhancing}
M.\ Safari, S.\ Movahedi, A.\ Moradi, and A.\ Zhu.
\newblock Enhancing DNA foundation models to address masking inefficiencies.
\newblock \textit{arXiv:2502.18405}, 2025.

\bibitem{ji2021dna}
Y.\ Ji, Z.\ Zhou, H.\ Liu, and R.\ V.\ Davuluri.
\newblock DNABERT: Pre-trained bidirectional encoder representations from
transformers model for DNA-language in genome.
\newblock \textit{Bioinformatics}, 37(15):2112--2120, 2021.

\bibitem{luo2024gue}
Y.\ Luo, B.\ Bhardwaj, K.\ Sundaram, and P.\ Huang.
\newblock GUE: Genomic understanding evaluation.
\newblock \textit{bioRxiv}, 2024.

\bibitem{lieber2024jamba}
O.\ Lieber, B.\ Lenz, H.\ Bata, et al.
\newblock Jamba: A hybrid Transformer-Mamba language model.
\newblock \textit{arXiv:2403.19887}, 2024.

\bibitem{ren2024samba}
L.\ Ren, Y.\ Liu, Y.\ Lu, Y.\ Shen, C.\ Liang, and W.\ Chen.
\newblock Samba: Simple hybrid state space models for efficient unlimited
context language modeling.
\newblock \textit{arXiv:2406.07522}, 2024.

\bibitem{hatamizadeh2025mamba}
A.\ Hatamizadeh and J.\ Kautz.
\newblock MambaVision: A hybrid Mamba-Transformer vision backbone.
\newblock \textit{CVPR}, 2025.

\bibitem{hopfieldlayers}
H.\ Ramsauer et al.
\newblock Hopfield layers.
\newblock \url{https://github.com/ml-jku/hopfield-layers}, 2021.

\bibitem{mambassm}
A.\ Gu and T.\ Dao.
\newblock Mamba: Linear-time sequence modeling with selective state spaces.
\newblock \url{https://github.com/state-spaces/mamba}, 2023.

\end{thebibliography}

\end{document}
